\documentclass[10pt,a4paper]{amsart}
\usepackage{amsmath,amssymb,amsthm,graphicx}
% Retained minimal preamble for the Stacks Project chapter varieties.tex.
% Verbatim excerpt of the chapter preamble: Xy-pic setup (native preamble lines 14-18), the native theorem
% declarations with the mathematical macros (native preamble lines 154-217), and the project own no-op
% definitions of the reference/slogan/history display environments (native wrapper.tex lines 1-3) that the
% counted statement uses for its slogan. Copied mechanically; nothing added, removed or redefined.
% For commutative diagrams we use Xy-pic
\usepackage[all]{xy}

% We use 2cell for 2-commutative diagrams.
\xyoption{2cell}
% Theorem environments.
%
\theoremstyle{plain}
\newtheorem{theorem}[subsection]{Theorem}
\newtheorem{proposition}[subsection]{Proposition}
\newtheorem{lemma}[subsection]{Lemma}

\theoremstyle{definition}
\newtheorem{definition}[subsection]{Definition}
\newtheorem{example}[subsection]{Example}
\newtheorem{exercise}[subsection]{Exercise}
\newtheorem{situation}[subsection]{Situation}

\theoremstyle{remark}
\newtheorem{remark}[subsection]{Remark}
\newtheorem{remarks}[subsection]{Remarks}

\numberwithin{equation}{subsection}

% Macros
%
\def\lim{\mathop{\mathrm{lim}}\nolimits}
\def\colim{\mathop{\mathrm{colim}}\nolimits}
\def\Spec{\mathop{\mathrm{Spec}}}
\def\Hom{\mathop{\mathrm{Hom}}\nolimits}
\def\Ext{\mathop{\mathrm{Ext}}\nolimits}
\def\SheafHom{\mathop{\mathcal{H}\!\mathit{om}}\nolimits}
\def\SheafExt{\mathop{\mathcal{E}\!\mathit{xt}}\nolimits}
\def\Sch{\mathit{Sch}}
\def\Mor{\mathop{\mathrm{Mor}}\nolimits}
\def\Ob{\mathop{\mathrm{Ob}}\nolimits}
\def\Sh{\mathop{\mathit{Sh}}\nolimits}
\def\NL{\mathop{N\!L}\nolimits}
\def\CH{\mathop{\mathrm{CH}}\nolimits}
\def\proetale{{pro\text{-}\acute{e}tale}}
\def\etale{{\acute{e}tale}}
\def\QCoh{\mathit{QCoh}}
\def\Ker{\mathop{\mathrm{Ker}}}
\def\Im{\mathop{\mathrm{Im}}}
\def\Coker{\mathop{\mathrm{Coker}}}
\def\Coim{\mathop{\mathrm{Coim}}}

% Boxtimes
%
\DeclareMathSymbol{\boxtimes}{\mathbin}{AMSa}{"02}

%
% Macros for moduli stacks/spaces
%
\def\QCohstack{\mathcal{QC}\!\mathit{oh}}
\def\Cohstack{\mathcal{C}\!\mathit{oh}}
\def\Spacesstack{\mathcal{S}\!\mathit{paces}}
\def\Quotfunctor{\mathrm{Quot}}
\def\Hilbfunctor{\mathrm{Hilb}}
\def\Curvesstack{\mathcal{C}\!\mathit{urves}}
\def\Polarizedstack{\mathcal{P}\!\mathit{olarized}}
\def\Complexesstack{\mathcal{C}\!\mathit{omplexes}}
% \Pic is the operator that assigns to X its picard group, usage \Pic(X)
% \Picardstack_{X/B} denotes the Picard stack of X over B
% \Picardfunctor_{X/B} denotes the Picard functor of X over B
\def\Pic{\mathop{\mathrm{Pic}}\nolimits}
\def\Picardstack{\mathcal{P}\!\mathit{ic}}
\def\Picardfunctor{\mathrm{Pic}}
\def\Deformationcategory{\mathcal{D}\!\mathit{ef}}
\newenvironment{reference}{}{}
\newenvironment{slogan}{}{}
\newenvironment{history}{}{}
\usepackage{cleveref}
\crefname{originalreference}{original reference}{original references}
\Crefname{originalreference}{Original reference}{Original references}
\makeatletter
\newlabel{morphisms-definition-normalization}{{035N}{source}{}{sourcepointer-1}{}}
\newlabel{morphisms-definition-normalization@cref}{{[originalreference][0][]035N}{[1][0][]source}{}{}{}}
\makeatother
\crefname{originalreference}{original reference}{original references}
\Crefname{originalreference}{Original reference}{Original references}
\makeatletter
\newlabel{morphisms-lemma-stalk-normalization}{{0C3B}{source}{}{sourcepointer-2}{}}
\newlabel{morphisms-lemma-stalk-normalization@cref}{{[originalreference][0][]0C3B}{[1][0][]source}{}{}{}}
\makeatother
\crefname{originalreference}{original reference}{original references}
\Crefname{originalreference}{Original reference}{Original references}
\makeatletter
\newlabel{algebra-lemma-integral-no-inclusion}{{00GT}{source}{}{sourcepointer-3}{}}
\newlabel{algebra-lemma-integral-no-inclusion@cref}{{[originalreference][0][]00GT}{[1][0][]source}{}{}{}}
\makeatother
\crefname{originalreference}{original reference}{original references}
\Crefname{originalreference}{Original reference}{Original references}
\makeatletter
\newlabel{algebra-lemma-integral-going-up}{{00GU}{source}{}{sourcepointer-4}{}}
\newlabel{algebra-lemma-integral-going-up@cref}{{[originalreference][0][]00GU}{[1][0][]source}{}{}{}}
\makeatother
\crefname{originalreference}{original reference}{original references}
\Crefname{originalreference}{Original reference}{Original references}
\makeatletter
\newlabel{more-algebra-lemma-number-of-branches-1}{{0C37}{source}{}{sourcepointer-5}{}}
\newlabel{more-algebra-lemma-number-of-branches-1@cref}{{[originalreference][0][]0C37}{[1][0][]source}{}{}{}}
\makeatother
\title{Branches and geometric branches at a point are counted by the fibres of the normalization}
\author{The Stacks Project Authors}
\date{stacks-project:varieties.tex:local-snapshot-20261011; GNU FDL 1.2 or later}
\begin{document}
\maketitle
\noindent\textit{Editorial note.} All delivered mathematical text is the author's own, copied verbatim from the pinned Stacks Project chapter file varieties.tex; no statement or proof step was rewritten, reordered or removed.\par\medskip
\noindent\textit{Editorial note.} The counted span cites five cross-chapter labels (including a definition), each recorded in assembly.external\_references with its canonical tag as a navigation pointer.\par\medskip
\noindent\textit{Editorial note.} The proof cites no bibliography item inside the counted span.\par\medskip
\noindent\textit{Allowed standard prerequisites (editorial).}\begin{itemize}
\item Stalk of the normalization as an integral closure. Stacks Project, Morphisms, Lemma morphisms-lemma-stalk-normalization (tag 035Q). Identifies the stalk of the pushforward of the structure sheaf along the normalization with the integral closure of O\_\{X,x\}.
\item Integral extensions and the correspondence of primes; algebraic branch count. Stacks Project, Algebra, Lemmas algebra-lemma-integral-no-inclusion (tag 00GS) and algebra-lemma-integral-going-up (tag 00GT); More on Algebra, Lemma more-algebra-lemma-number-of-branches-1 (tag 0C1R). The first two turn the fibre over x into maximal ideals of the integral closure and the third supplies the branch formulas.
\end{itemize}\medskip
\ifdefined\hypertarget\hypertarget{sourcepointer-1}{}\fi
\noindent\textit{Source reference.} Original reference 035N: sources/stacks-chow/tags.txt, source lines 4134--4134. Cross-chapter standard prerequisite cited by the author inside the counted span; it is used as a pointer and its own proof is not part of this item.\par\medskip
\ifdefined\hypertarget\hypertarget{sourcepointer-2}{}\fi
\noindent\textit{Source reference.} Original reference 0C3B: sources/stacks-chow/tags.txt, source lines 15186--15186. Cross-chapter standard prerequisite cited by the author inside the counted span; it is used as a pointer and its own proof is not part of this item.\par\medskip
\ifdefined\hypertarget\hypertarget{sourcepointer-3}{}\fi
\noindent\textit{Source reference.} Original reference 00GT: sources/stacks-chow/tags.txt, source lines 619--619. Cross-chapter standard prerequisite cited by the author inside the counted span; it is used as a pointer and its own proof is not part of this item.\par\medskip
\ifdefined\hypertarget\hypertarget{sourcepointer-4}{}\fi
\noindent\textit{Source reference.} Original reference 00GU: sources/stacks-chow/tags.txt, source lines 620--620. Cross-chapter standard prerequisite cited by the author inside the counted span; it is used as a pointer and its own proof is not part of this item.\par\medskip
\ifdefined\hypertarget\hypertarget{sourcepointer-5}{}\fi
\noindent\textit{Source reference.} Original reference 0C37: sources/stacks-chow/tags.txt, source lines 15182--15182. Cross-chapter standard prerequisite cited by the author inside the counted span; it is used as a pointer and its own proof is not part of this item.\par\medskip
\begin{lemma}
\label{lemma-number-of-branches}
Let $X$ be a scheme. Assume every quasi-compact open of $X$ has
finitely many irreducible components. Let $\nu : X^\nu \to X$
be the normalization of $X$. Let $x \in X$.
\begin{enumerate}
\item The number of branches of $X$ at $x$ is the number of
inverse images of $x$ in $X^\nu$.
\item The number of geometric branches of $X$ at $x$ is
$\sum_{\nu(x^\nu) = x} [\kappa(x^\nu) : \kappa(x)]_s$.
\end{enumerate}
\end{lemma}

\begin{proof}
First note that the assumption on $X$ exactly means that the
normalization is defined, see
Morphisms, Definition \ref{morphisms-definition-normalization}.
Then the stalk $A' = (\nu_*\mathcal{O}_{X^\nu})_x$ is the
integral closure of $A = \mathcal{O}_{X, x}$ in the total
ring of fractions of $A_{red}$, see
Morphisms, Lemma \ref{morphisms-lemma-stalk-normalization}.
Since $\nu$ is an integral morphism, we see that
the points of $X^\nu$ lying over $x$ correspond
to the primes of $A'$ lying over the maximal ideal $\mathfrak m$
of $A$. As $A \to A'$ is integral, this is the same thing as
the maximal ideals of $A'$
(Algebra, Lemmas \ref{algebra-lemma-integral-no-inclusion} and
\ref{algebra-lemma-integral-going-up}).
Thus the lemma now follows
from its algebraic counterpart:
More on Algebra, Lemma \ref{more-algebra-lemma-number-of-branches-1}.
\end{proof}

\end{document}
